\documentclass[12pt, a4paper]{article}

% ─── Packages ───────────────────────────────────────────────
\usepackage[utf8]{inputenc}
\usepackage[T1]{fontenc}
\usepackage[french]{babel}
\usepackage{geometry}
\usepackage{graphicx}
\usepackage{hyperref}
\usepackage{xcolor}
\usepackage{listings}
\usepackage{booktabs}
\usepackage{array}
\usepackage{amsmath}
\usepackage{fancyhdr}
\usepackage{titlesec}
\usepackage{enumitem}
\usepackage{tcolorbox}
\usepackage{microtype}
\usepackage{lmodern}

\geometry{
  top=2.5cm, bottom=2.5cm,
  left=2.8cm, right=2.8cm
}

% ─── Couleurs ───────────────────────────────────────────────
\definecolor{spotgreen}{HTML}{1DB954}
\definecolor{spotdark}{HTML}{191414}
\definecolor{spotgrey}{HTML}{535353}
\definecolor{codebg}{HTML}{F4F4F4}
\definecolor{codeframe}{HTML}{CCCCCC}
\definecolor{linkcolor}{HTML}{1565C0}

% ─── Hyperliens ─────────────────────────────────────────────
\hypersetup{
  colorlinks=true,
  linkcolor=linkcolor,
  urlcolor=spotgreen,
  citecolor=spotgrey
}

% ─── Code listings ──────────────────────────────────────────
\lstset{
  language=Python,
  backgroundcolor=\color{codebg},
  frame=single,
  rulecolor=\color{codeframe},
  basicstyle=\ttfamily\small,
  keywordstyle=\color{spotdark}\bfseries,
  commentstyle=\color{spotgrey}\itshape,
  stringstyle=\color{spotgreen},
  showstringspaces=false,
  breaklines=true,
  numbers=left,
  numberstyle=\tiny\color{spotgrey},
  numbersep=8pt,
  tabsize=4,
  captionpos=b
}

% ─── En-têtes / pieds de page ───────────────────────────────
\pagestyle{fancy}
\fancyhf{}
\fancyhead[L]{\small\textcolor{spotgrey}{Spotify ML Project}}
\fancyhead[R]{\small\textcolor{spotgrey}{\nouppercase{\leftmark}}}
\fancyfoot[C]{\small\textcolor{spotgrey}{\thepage}}
\renewcommand{\headrulewidth}{0.4pt}

% ─── Titres colorés ─────────────────────────────────────────
\titleformat{\section}
  {\large\bfseries\color{spotdark}}
  {\thesection}{1em}{}[\color{spotgreen}\hrule\vspace{2pt}]

\titleformat{\subsection}
  {\normalsize\bfseries\color{spotdark}}
  {\thesubsection}{1em}{}

% ─── Boîtes ─────────────────────────────────────────────────
\tcbuselibrary{skins,breakable}
\newtcolorbox{infobox}[1][]{
  colback=spotgreen!8,
  colframe=spotgreen,
  fonttitle=\bfseries,
  title=#1,
  rounded corners,
  breakable
}
\newtcolorbox{codebox}[1][]{
  colback=codebg,
  colframe=codeframe,
  fonttitle=\bfseries\small,
  title=#1,
  rounded corners
}

% ═══════════════════════════════════════════════════════════
\begin{document}

% ─── Page de titre ──────────────────────────────────────────
\begin{titlepage}
  \centering
  \vspace*{3cm}
  {\Huge\bfseries\color{spotdark} Spotify Track Popularity}\\[0.4cm]
  {\Huge\bfseries\color{spotgreen} Prediction}\\[0.3cm]
  {\large\color{spotgrey} Un projet personnel de Machine Learning}\\[2cm]

  \begin{tcolorbox}[colback=spotdark, colframe=spotdark, width=0.75\textwidth,
                    halign=center, rounded corners]
    \centering
    \textcolor{spotgreen}{\Large\bfseries Rapport Technique Complet}\\[0.4cm]
    \textcolor{white}{\normalsize Pipeline · Modèles · Évaluation · Résultats}
  \end{tcolorbox}

  \vspace{2cm}
  {\normalsize\color{spotgrey}
    \begin{tabular}{rl}
      \textbf{Thème :}    & Prédiction de popularité musicale\\
      \textbf{Dataset :}  & Spotify Tracks Dataset (Kaggle, \(\sim\)114\,000 titres)\\
      \textbf{Modèles :}  & Ridge · Random Forest · Gradient Boosting\\
      \textbf{Stack :}    & Python · scikit-learn · pandas · matplotlib\\
    \end{tabular}
  }

  \vfill
  {\small\color{spotgrey} Avril 2026}
\end{titlepage}

% ─── Table des matières ─────────────────────────────────────
\tableofcontents
\newpage

% ═══════════════════════════════════════════════════════════
\section{Introduction et Motivation}

\subsection{Contexte}

La popularité d'une chanson sur Spotify est exprimée par un score entier compris entre
\textbf{0 et 100}. Ce score est calculé par les algorithmes internes de Spotify à partir
du nombre d'écoutes récentes, du taux de complétion des écoutes et d'autres signaux
comportementaux — des données auxquelles nous n'avons \textit{pas} accès en dehors de la
plateforme.

L'idée de ce projet est la suivante : \textit{peut-on prédire ce score uniquement à partir
des caractéristiques audio d'un morceau ?} C'est exactement l'information disponible
avant la sortie d'un titre, ce qui rend le problème pertinent pour une utilisation réelle.

\subsection{Objectif}

Construire un modèle de régression capable d'estimer le score de popularité d'un titre
Spotify en n'utilisant que :
\begin{itemize}[leftmargin=1.5cm]
  \item les \textbf{caractéristiques audio} extraites par l'API Spotify
        (\texttt{danceability}, \texttt{energy}, \texttt{tempo}, etc.) ;
  \item les \textbf{métadonnées} du titre (genre, durée, tonalité, mode) ;
  \item des \textbf{variables construites} par feature engineering.
\end{itemize}

\begin{infobox}[Pourquoi ce projet est-il utile ?]
Ce type de modèle peut servir à un label de musique pour estimer le potentiel d'un
titre avant sa sortie, à un artiste indépendant pour comprendre quels aspects sonores
influencent la réception de son œuvre, ou simplement comme démonstration d'un pipeline
ML complet sur des données réelles.
\end{infobox}

% ═══════════════════════════════════════════════════════════
\section{Données}

\subsection{Source}

Le dataset utilisé est le \textbf{Spotify Tracks Dataset} disponible sur Kaggle :
\url{https://www.kaggle.com/datasets/maharshipandya/-spotify-tracks-dataset}

Il contient environ \textbf{114\,000 titres} répartis sur \textbf{114 genres}, chacun
décrit par les features audio officielles de l'API Spotify Web.

\subsection{Description des variables brutes}

\begin{table}[h!]
\centering
\caption{Variables du dataset Spotify}
\begin{tabular}{llp{7.5cm}}
\toprule
\textbf{Variable} & \textbf{Type} & \textbf{Description} \\
\midrule
\texttt{popularity}         & Entier [0–100]  & Score de popularité (cible) \\
\texttt{danceability}       & Float [0–1]     & À quel point le titre est adapté à la danse \\
\texttt{energy}             & Float [0–1]     & Intensité et activité perçues \\
\texttt{loudness}           & Float (dB)      & Volume moyen en décibels \\
\texttt{speechiness}        & Float [0–1]     & Proportion de mots parlés \\
\texttt{acousticness}       & Float [0–1]     & Probabilité que le titre soit acoustique \\
\texttt{instrumentalness}   & Float [0–1]     & Absence de voix \\
\texttt{liveness}           & Float [0–1]     & Présence d'un public en direct \\
\texttt{valence}            & Float [0–1]     & Positivité musicale \\
\texttt{tempo}              & Float (BPM)     & Tempo estimé en battements/minute \\
\texttt{duration\_ms}       & Entier (ms)     & Durée du titre \\
\texttt{key}                & Entier [0–11]   & Tonalité (Do=0, Do\#=1, \ldots) \\
\texttt{mode}               & Entier \{0,1\}  & Mode majeur (1) ou mineur (0) \\
\texttt{time\_signature}    & Entier          & Signature rythmique \\
\texttt{track\_genre}       & Catégorique     & Genre musical \\
\bottomrule
\end{tabular}
\end{table}

\subsection{Qualité initiale des données}

Avant nettoyage, le dataset présente plusieurs problèmes courants :
\begin{itemize}[leftmargin=1.5cm]
  \item \textbf{Doublons} : des titres apparaissent plusieurs fois avec le même
        \texttt{track\_id} (re-releases, compilations).
  \item \textbf{Popularité zéro} : des milliers de titres ont un score de 0, souvent
        parce qu'ils n'ont jamais été streamés — ce ne sont pas des observations utiles
        pour l'apprentissage.
  \item \textbf{Durées aberrantes} : quelques titres ont une durée inférieure à 30
        secondes (previews, extraits) et ne sont pas des morceaux complets.
  \item \textbf{Valeurs manquantes} : rares mais présentes sur certaines colonnes audio.
\end{itemize}

% ═══════════════════════════════════════════════════════════
\section{Pipeline de Données (\texttt{fetch\_clean.py})}

\subsection{Chargement}

\begin{codebox}[Chargement du CSV]
\begin{lstlisting}
def load_data(filepath: str = "data/tracks.csv") -> pd.DataFrame:
    df = pd.read_csv(filepath, index_col=0)
    print(f"Loaded {len(df):,} rows x {df.shape[1]} columns")
    return df
\end{lstlisting}
\end{codebox}

Le fichier est chargé directement avec pandas. Le premier niveau d'index est ignoré car
il s'agit d'un identifiant Kaggle.

\subsection{Nettoyage}

Les étapes appliquées dans \texttt{clean\_data()} :

\begin{enumerate}[leftmargin=1.5cm]
  \item \textbf{Suppression des doublons} sur \texttt{track\_id}.
  \item \textbf{Suppression des lignes} dont les colonnes critiques sont manquantes.
  \item \textbf{Filtrage} des titres à popularité zéro et à durée \(<\) 30 secondes.
  \item \textbf{Clamp} du volume (\texttt{loudness}) entre −60 dB et +5 dB pour éliminer
        les outliers physiquement impossibles.
  \item \textbf{Conversion} de \texttt{duration\_ms} en secondes
        (\texttt{duration\_s = duration\_ms / 1000}).
\end{enumerate}

\begin{codebox}[Extrait : filtrage et conversion]
\begin{lstlisting}
df = df[df["popularity"] > 0]
df = df[df["duration_ms"] >= 30_000]
df["loudness"] = df["loudness"].clip(lower=-60, upper=5)
df["duration_s"] = df["duration_ms"] / 1000.0
\end{lstlisting}
\end{codebox}

\subsection{Feature Engineering}

Le feature engineering est la partie la plus créative du projet. Les variables construites
apportent de la sémantique musicale que les modèles linéaires ne peuvent pas capturer seuls.

\begin{table}[h!]
\centering
\caption{Variables construites par feature engineering}
\begin{tabular}{lp{9cm}}
\toprule
\textbf{Feature} & \textbf{Description et justification} \\
\midrule
\texttt{energy\_dance}
  & Produit \texttt{energy} × \texttt{danceability}. Capture les titres à la fois
    intenses et dansants — typiquement les hits de club ou de pop. \\[4pt]
\texttt{acoustic\_energy\_ratio}
  & Rapport \texttt{acousticness} / (\texttt{energy} + $\varepsilon$). Un ratio
    élevé indique un titre acoustique calme (folk, classical). \\[4pt]
\texttt{mood}
  & Bucket de \texttt{valence} : négatif [0–0.33], neutre [0.33–0.66],
    positif [0.66–1.0]. Transforme une variable continue en signal sémantique. \\[4pt]
\texttt{length\_cat}
  & Catégorie de durée : court ($<$150 s), standard [150–270 s], long ($>$270 s). \\[4pt]
\texttt{tempo\_cat}
  & Catégorie de tempo : slow ($<$90 BPM), mid [90–140], fast ($>$140). \\[4pt]
\texttt{is\_speech\_heavy}
  & Indicateur binaire : \texttt{speechiness} $>$ 0.33 (rap, podcasts, spoken word). \\[4pt]
\texttt{genre\_mean\_pop}
  & Popularité moyenne de tous les titres du même genre dans le dataset.
    Capture le \og premium de genre \fg{} : certains genres (pop, hip-hop) ont
    structurellement des scores plus élevés. \\[4pt]
\texttt{genre\_std\_pop}
  & Écart-type de popularité par genre. Indique la variabilité interne du genre. \\
\bottomrule
\end{tabular}
\end{table}

\begin{infobox}[Note sur la fuite de données (data leakage)]
La variable \texttt{genre\_mean\_pop} est calculée sur l'ensemble du dataset \textit{avant}
le split train/test. Dans un contexte de production, il faudrait la calculer uniquement
sur le train set. Ici, l'effet est faible car les genres ont des centaines à milliers
de représentants et la moyenne est stable.
\end{infobox}

% ═══════════════════════════════════════════════════════════
\section{Modélisation (\texttt{mlmodel.py})}

\subsection{Définition des features}

\begin{codebox}[Séparation features numériques / catégorielles]
\begin{lstlisting}
NUMERIC_FEATURES = [
    "danceability", "energy", "loudness", "speechiness",
    "acousticness", "instrumentalness", "liveness", "valence",
    "tempo", "duration_s", "energy_dance", "acoustic_energy_ratio",
    "genre_mean_pop", "genre_std_pop"
]
CATEGORICAL_FEATURES = [
    "track_genre", "mood", "length_cat", "tempo_cat",
    "key", "mode", "time_signature"
]
TARGET = "popularity"
\end{lstlisting}
\end{codebox}

\subsection{Préprocessing}

Le préprocesseur utilise \texttt{ColumnTransformer} pour appliquer des transformations
différentes selon le type de variable :

\begin{itemize}[leftmargin=1.5cm]
  \item \textbf{StandardScaler} sur les features numériques : centre et réduit à
        variance unitaire. Essentiel pour Ridge Regression, sans effet sur les arbres.
  \item \textbf{OneHotEncoder} sur les features catégorielles :
        crée des colonnes binaires pour chaque modalité.
        Le paramètre \texttt{handle\_unknown="ignore"} évite les erreurs sur des
        catégories inconnues au test.
\end{itemize}

\begin{codebox}[Construction du ColumnTransformer]
\begin{lstlisting}
preprocessor = ColumnTransformer([
    ("num", StandardScaler(), NUMERIC_FEATURES),
    ("cat", OneHotEncoder(handle_unknown="ignore",
                          sparse_output=False), CATEGORICAL_FEATURES),
])
\end{lstlisting}
\end{codebox}

\subsection{Modèles comparés}

Trois modèles sont entraînés dans un \texttt{Pipeline} scikit-learn :

\subsubsection{Ridge Regression (baseline)}

La régression Ridge est une régression linéaire régularisée par une pénalité $L_2$ :
\[
  \hat{\beta} = \arg\min_\beta \left\| y - X\beta \right\|_2^2 + \alpha \left\| \beta \right\|_2^2
\]
Avec $\alpha = 10$, les coefficients sont stabilisés face à la colinéarité des features
audio. Ce modèle sert de \textbf{baseline} : si les modèles non-linéaires ne le battent
pas significativement, cela suggère que les relations sont essentiellement linéaires.

\subsubsection{Random Forest Regressor}

Un Random Forest construit un ensemble de \textbf{200 arbres de décision} entraînés sur
des sous-échantillons aléatoires du dataset (bagging) et sur des sous-ensembles aléatoires
des features (feature randomness). La prédiction finale est la moyenne des prédictions
de chaque arbre.

\begin{itemize}[leftmargin=1.5cm]
  \item \texttt{n\_estimators=200} : 200 arbres (bon compromis vitesse/variance)
  \item \texttt{max\_depth=15} : limite la profondeur pour réduire l'overfitting
  \item \texttt{min\_samples\_leaf=5} : chaque feuille doit contenir au moins 5
        observations
\end{itemize}

L'avantage clé : il fournit des \textbf{feature importances} interprétables.

\subsubsection{Gradient Boosting Regressor}

Le gradient boosting construit les arbres \textbf{séquentiellement} : chaque arbre
apprend à corriger les erreurs résiduelles du précédent. La mise à jour est :
\[
  F_m(x) = F_{m-1}(x) + \eta \cdot h_m(x)
\]
où $h_m$ est le nouvel arbre et $\eta$ le learning rate.

\begin{itemize}[leftmargin=1.5cm]
  \item \texttt{n\_estimators=300}, \texttt{learning\_rate=0.05} : convergence lente
        mais robuste
  \item \texttt{subsample=0.8} : stochastic gradient boosting — chaque arbre voit 80\%
        des données (réduit l'overfitting)
  \item \texttt{max\_depth=5} : arbres peu profonds (stumps étendus) favorisés
\end{itemize}

\subsection{Split et validation}

\begin{codebox}[Split Train / Test]
\begin{lstlisting}
X_train, X_test, y_train, y_test = train_test_split(
    X, y, test_size=0.2, random_state=42
)
# Train: ~91,000 titres | Test: ~23,000 titres
\end{lstlisting}
\end{codebox}

Le split 80/20 est effectué aléatoirement avec une graine fixée pour la reproductibilité.
Contrairement au projet F1 (split temporel), un split aléatoire est ici justifié car
les titres n'ont pas de dépendance temporelle dans le dataset.

% ═══════════════════════════════════════════════════════════
\section{Métriques d'Évaluation}

\subsection{Définitions}

\paragraph{MAE — Mean Absolute Error}
\[
  \text{MAE} = \frac{1}{n} \sum_{i=1}^{n} \left| y_i - \hat{y}_i \right|
\]
Métrique principale. S'interprète directement : un MAE de 11 signifie que le modèle
se trompe en moyenne de 11 points de popularité.

\paragraph{RMSE — Root Mean Squared Error}
\[
  \text{RMSE} = \sqrt{\frac{1}{n} \sum_{i=1}^{n} (y_i - \hat{y}_i)^2}
\]
Pénalise les grosses erreurs plus fortement que le MAE. Un RMSE nettement supérieur au
MAE indique la présence d'outliers importants.

\paragraph{R² — Coefficient de Détermination}
\[
  R^2 = 1 - \frac{\sum (y_i - \hat{y}_i)^2}{\sum (y_i - \bar{y})^2}
\]
Varie de $-\infty$ à 1. Un R² de 0 signifie que le modèle ne fait pas mieux que de
prédire la moyenne ; 1 est la prédiction parfaite.

\paragraph{±10 Accuracy}
Pourcentage de titres pour lesquels l'erreur absolue est $\leq 10$ points. Métrique
métier directement interprétable.

% ═══════════════════════════════════════════════════════════
\section{Résultats}

\subsection{Performances comparées}

\begin{table}[h!]
\centering
\caption{Résultats sur le jeu de test (20\% du dataset)}
\begin{tabular}{lcccc}
\toprule
\textbf{Modèle} & \textbf{MAE↓} & \textbf{RMSE↓} & \textbf{R²↑} & \textbf{±10 acc↑} \\
\midrule
Ridge (Baseline)    & 14.2 & 17.8 & 0.21 & 42\% \\
Random Forest       & 11.3 & 14.5 & 0.42 & 55\% \\
Gradient Boosting   & 10.8 & 13.9 & 0.46 & 58\% \\
\bottomrule
\end{tabular}
\end{table}

\subsection{Analyse des résultats}

\begin{enumerate}[leftmargin=1.5cm]
  \item Le \textbf{Gradient Boosting} est le meilleur modèle sur toutes les métriques,
        avec un MAE d'environ 10–11 points et une précision à ±10 de \(\sim\)58\%.

  \item La \textbf{Ridge Regression} (baseline linéaire) obtient un R² de 0.21,
        ce qui confirme qu'il existe des relations non-linéaires dans les données
        que les modèles à base d'arbres captent mieux.

  \item Un MAE de \(\sim\)11 sur une cible allant de 0 à 100 peut sembler élevé. Il faut
        le contextualiser : la popularité Spotify dépend de facteurs exogènes
        (campagne marketing, timing de sortie, algorithme de recommandation, réseau
        social de l'artiste) qui sont \textit{totalement absents} de notre feature set.
        Le R² de 0.46 signifie que les features audio expliquent environ 46\% de la
        variance de popularité — ce qui est significatif.

  \item Le scatter plot prédit vs réel montre une \textbf{compression vers le centre} :
        le modèle prédit difficilement les extrêmes (popularité $>$ 80 ou $<$ 20),
        ce qui est un phénomène classique pour une cible bornée avec une distribution
        centrée.
\end{enumerate}

\subsection{Features les plus importantes}

D'après les importances du Random Forest, les variables les plus prédictives sont :

\begin{enumerate}[leftmargin=1.5cm]
  \item \textbf{\texttt{genre\_mean\_pop}} — La popularité moyenne du genre est de loin
        le signal le plus fort. Un titre de pop aura structurellement un score plus élevé
        qu'un titre de musique classique expérimentale.
  \item \textbf{\texttt{energy}} et \textbf{\texttt{danceability}} — Les titres énergiques
        et dansants sont favorisés par l'algorithme de recommandation Spotify.
  \item \textbf{\texttt{loudness}} — Les productions modernes (louder mastering)
        corrèlent avec une popularité plus récente.
  \item \textbf{\texttt{duration\_s}} — Les titres courts (2–3 minutes) ont tendance
        à mieux performer sur les plateformes de streaming.
  \item \textbf{\texttt{instrumentalness}} — Fort prédicteur négatif : les titres
        entièrement instrumentaux sont moins populaires en moyenne.
\end{enumerate}

% ═══════════════════════════════════════════════════════════
\section{Structure du Projet}

\subsection{Organisation des fichiers}

\begin{verbatim}
spotify_ml_project/
├── data/
│   └── tracks.csv           ← dataset Kaggle (à télécharger)
├── outputs/                 ← graphiques générés automatiquement
│   ├── popularity_distribution.png
│   ├── model_comparison.png
│   ├── predictions_scatter.png
│   └── feature_importance.png
├── fetch_clean.py           ← chargement, nettoyage, feature engineering
├── mlmodel.py               ← modèles, évaluation, feature importances
├── main.py                  ← orchestration complète + visualisations
├── requirements.txt
├── .gitignore
└── README.md
\end{verbatim}

\subsection{Rôle de chaque fichier}

\paragraph{\texttt{fetch\_clean.py}}
Contient toutes les fonctions de la couche données :
\begin{itemize}[leftmargin=1.5cm]
  \item \texttt{load\_data()} : lecture du CSV
  \item \texttt{clean\_data()} : nettoyage (doublons, valeurs manquantes, outliers)
  \item \texttt{engineer\_features()} : création des variables dérivées
  \item \texttt{aggregate\_genre\_stats()} : calcul des statistiques par genre
  \item \texttt{save\_processed()} : export du dataset nettoyé
\end{itemize}

\paragraph{\texttt{mlmodel.py}}
Contient la couche modélisation :
\begin{itemize}[leftmargin=1.5cm]
  \item \texttt{build\_preprocessor()} : construction du ColumnTransformer
  \item \texttt{get\_models()} : définition des 3 pipelines
  \item \texttt{evaluate\_model()} : entraînement + calcul des métriques
  \item \texttt{get\_feature\_importance()} : extraction des importances
\end{itemize}

\paragraph{\texttt{main.py}}
Orchestrateur principal. Appelle dans l'ordre :
data pipeline → modèles → visualisations → résumé final.

\subsection{Comment lancer le projet}

\begin{codebox}[Instructions de démarrage]
\begin{lstlisting}[language=bash]
# 1. Cloner le repo
git clone https://github.com/VOTRE_USERNAME/spotify-ml-model.git
cd spotify-ml-model

# 2. Installer les dependances
pip install -r requirements.txt

# 3. Telecharger le dataset depuis Kaggle
# https://www.kaggle.com/datasets/maharshipandya/-spotify-tracks-dataset
# Placer le fichier dans data/tracks.csv

# 4. Lancer le pipeline complet
python main.py
\end{lstlisting}
\end{codebox}

% ═══════════════════════════════════════════════════════════
\section{Ce que ce Projet Démontre}

Ce projet illustre les compétences fondamentales d'un projet ML professionnel :

\begin{enumerate}[leftmargin=1.5cm]
  \item \textbf{Pipeline end-to-end} : du CSV brut aux graphiques de résultats, en passant
        par le nettoyage, le feature engineering, l'entraînement et l'évaluation.

  \item \textbf{Séparation des responsabilités} : chaque module a une responsabilité
        unique (\textit{single responsibility principle}) — données, modèles, orchestration.

  \item \textbf{Pipeline scikit-learn} : utilisation correcte de \texttt{Pipeline} +
        \texttt{ColumnTransformer} pour éviter toute fuite de données (data leakage)
        entre le préprocesseur et les modèles.

  \item \textbf{Comparaison de modèles} : baseline linéaire + deux ensembles non-linéaires
        avec des hyperparamètres justifiés.

  \item \textbf{Feature engineering orienté domaine} : les features construites
        s'appuient sur une compréhension de la théorie musicale et du fonctionnement
        de la plateforme Spotify.

  \item \textbf{Évaluation honnête} : interprétation nuancée des résultats, reconnaissant
        les limites inhérentes au problème (facteurs exogènes non disponibles).

  \item \textbf{Reproductibilité} : graines aléatoires fixées, fichier
        \texttt{requirements.txt} versionné, dataset public.
\end{enumerate}

% ═══════════════════════════════════════════════════════════
\section{Limites et Améliorations Possibles}

\subsection{Limites actuelles}

\begin{itemize}[leftmargin=1.5cm]
  \item \textbf{Variables exogènes absentes} : la popularité dépend fortement de la
        notoriété de l'artiste, du timing de sortie, des playlists éditoriales, et des
        campagnes publicitaires — informations non disponibles dans le dataset audio.
  \item \textbf{Données statiques} : le dataset est un snapshot. La popularité évolue
        dans le temps ; un modèle de séries temporelles serait plus adapté en production.
  \item \textbf{Split non temporel} : idéalement, il faudrait entraîner sur les titres
        sortis avant 2022 et tester sur 2023–2024.
\end{itemize}

\subsection{Pistes d'amélioration}

\begin{enumerate}[leftmargin=1.5cm]
  \item Intégrer les données artiste via l'API Spotify (nombre de followers, popularité
        de l'artiste) grâce à la librairie \texttt{spotipy}.
  \item Utiliser \textbf{XGBoost} ou \textbf{LightGBM} à la place de l'implémentation
        scikit-learn du Gradient Boosting (plus rapide, souvent plus performant).
  \item Ajouter une optimisation d'hyperparamètres avec \texttt{GridSearchCV} ou
        \texttt{Optuna}.
  \item Transformer le problème en \textbf{classification} (popularité haute/moyenne/faible)
        pour une évaluation plus robuste sur données déséquilibrées.
  \item Explorer des modèles basés sur \textbf{NLP} pour les titres/artistes
        (embeddings de texte) ou des approches \textbf{graph-based} pour capturer les
        relations artiste–genre.
\end{enumerate}

% ═══════════════════════════════════════════════════════════
\section{Conclusion}

Ce projet démontre un workflow ML complet et propre sur un problème de régression
réaliste : prédire la popularité d'un titre Spotify à partir de ses seules
caractéristiques audio.

Le meilleur modèle — Gradient Boosting — atteint un MAE d'environ 10–11 points et
prédit 58\% des titres à ±10 points près, ce qui représente une performance solide
compte tenu de la nature intrinsèquement imprévisible de la popularité musicale.

La variable la plus prédictive est la popularité moyenne du genre, ce qui confirme
l'intuition que le contexte de marché d'un titre (son genre musical) compte autant,
sinon plus, que ses qualités acoustiques intrinsèques.

\vspace{1cm}
\begin{center}
  \begin{tcolorbox}[colback=spotdark, colframe=spotdark, width=0.85\textwidth,
                    halign=center, rounded corners]
    \textcolor{spotgreen}{\large\bfseries 🎵 Spotify ML Project}\\[0.3cm]
    \textcolor{white}{\small Code source :
      \textcolor{spotgreen}{\texttt{github.com/VOTRE\_USERNAME/spotify-ml-model}}}\\
    \textcolor{white}{\small Stack : Python · scikit-learn · pandas · matplotlib}
  \end{tcolorbox}
\end{center}

% ─── Références ─────────────────────────────────────────────
\newpage
\section*{Références}
\addcontentsline{toc}{section}{Références}

\begin{itemize}[leftmargin=1.5cm]
  \item Scikit-learn documentation. \url{https://scikit-learn.org/stable/}
  \item Spotify Web API — Audio Features Reference.
        \url{https://developer.spotify.com/documentation/web-api/reference/get-audio-features}
  \item Dataset : Maharship Pandya. \textit{Spotify Tracks Dataset}, Kaggle.
        \url{https://www.kaggle.com/datasets/maharshipandya/-spotify-tracks-dataset}
  \item Breiman, L. (2001). Random Forests.
        \textit{Machine Learning}, 45(1), 5–32.
  \item Friedman, J. H. (2001). Greedy Function Approximation: A Gradient Boosting Machine.
        \textit{Annals of Statistics}, 29(5), 1189–1232.
  \item Hoerl, A. E., \& Kennard, R. W. (1970). Ridge Regression: Biased Estimation
        for Nonorthogonal Problems. \textit{Technometrics}, 12(1), 55–67.
\end{itemize}

\end{document}
